\chapter{Neue Drucksachen}
\quellseite{10}

\section[Elfter Synodalbericht]{Elfter Synodalbericht des Kansas-Districts der deutschen ev.-luth. Synode von Missouri, Ohio u.\,a. St.}
St. Louis, Mo. 1903. Concordia Publishing House. 80 Seiten 9×6. Preis: 15 Cts.

Der Kansas-District unserer Synode hat auf seiner letzten Versammlung einen für alle Christen sehr wichtigen Gegenstand besprochen, nämlich das Kirchengehen. Diese lehrreichen Verhandlungen bilden den Haupttheil dieses Berichtes. Folgende sechs Punkte werden dabei ausgeführt: 1. Was unter Kirchengehen zu verstehen ist und wie die Christen es von Alters her geübt haben. 2. Ob solch Kirchengehen von Gott geboten sei, oder ob es einem Christen freistehe. 3. Welches der Nutzen und Segen des Kirchengehens sei und was für Gefahr und Schaden in der Versäumung des Kirchengehens liege. 4. Wie, 5. wo und 6. wann man zur Kirche gehen solle. Das alles wird des Weiteren schlicht und einfach, gründlich und deutlich dargelegt. Da nun immer noch manche Gemeinden über schlechten Kirchenbesuch klagen müssen, da in jeder Gemeinde sich immer noch solche finden, die im Kirchengehen nachlässig sind und so ihre Seele in Gefahr setzen, so wäre es wünschenswerth und würde auch gewiß nicht ohne Segen bleiben, wenn gerade auch dieser Bericht eine weite Verbreitung fände, von vielen gelesen und beherzigt würde. Sonst ist noch in dem Berichte zu erwähnen die Synodalrede des Präses und der Missionsbericht, der auch einen Ueberblick gibt über eins unserer wichtigsten Missionsfelder, über Oklahoma. G. M.

\section[Evangelisch-lutherischer Stadtmissionar]{Evangelisch-lutherischer Stadtmissionar}
Februarnummer. 4 Seiten 12×9. Preis: 50 Exemplare 20 Cts., 100 Exemplare 35 Cts. portofrei. Zu bestellen bei Rev. F. P. Merbitz, 3931 Dearborn St., Chicago, Ill.

Wir haben schon öfters an dieser Stelle den in Chicago erscheinenden, trefflich redigirten „Stadtmissionar“ als passenden Tractat zur Vertheilung an den Kirchthüren empfohlen. Daß wir diese neueste Nummer herausgreifen und ganz besonders besprechen, hat seinen guten Grund. Sie enthält nämlich eine am Sonntag nach dem schrecklichen Chicagoer Theaterbrand von einem dortigen treuverdienten Pastor gehaltene sehr zeitgemäße Predigt. Zum Text hat diese Predigt das Schriftwort 1 Sam. 3, 9.: „Rede, Herr, denn dein Knecht höret“ und zum Thema die Frage: „Was sagt uns Gott durch dieses schreckliche Ereigniß?“ Die Antwort ist eine doppelte: 1. „Gott hat Mißfallen am Theaterwesen“ und 2. „Gott will, daß wir unsere Seelen retten sollen.“ Im ersten Theil wird gezeigt, wie in dem heutigen Theater, der „Kirche der Welt“, alle Gebote Gottes mit Füßen getreten werden, wie es ein Ort ist, in dem beständig unnütze Narrentheidinge, unzüchtige Scherzreden, gottlose Spöttereien ausgesprochen und Zucht und Schamhaftigkeit auch in mancher andern Weise grob verletzt werden, namentlich durch die unanständige, anstößige Kleidung der Schauspielerinnen. Im zweiten Theil wird ausgeführt, wie Gott insbesondere der Jugend, den Jünglingen und Jungfrauen, zuruft: Eilet und errettet eure Seelen, stellet euch nicht dieser Welt gleich; wie er aber auch allen denen, die noch nicht wahrhaft bekehrt sind, und endlich auch denen, die schon im Glauben stehen, etwas zu sagen hat. Man kann die aus treuer, seelsorgerlicher Liebe hervorgegangene Predigt nicht ohne tiefe Bewegung lesen. Sie wird auch anderwärts, wo immer man gegen das heutige Theaterwesen zu zeugen Veranlassung hat, heilsamen Eindruck machen und unter Gottes Segen Frucht schaffen. Um Massenverbreitung zu erreichen, ist der Preis wirklich beispiellos billig angesetzt. L. F.

\section[Evangelisch-Lutherische Kirchenglocke]{„Evangelisch-Lutherische Kirchenglocke“}
Eine Mahn- und Weckstimme an alle Lutheraner. Herausgegeben von P. B. P. Rommensen, 1231 Kinnickinnic Ave., Milwaukee, Wis. 6 Jahrgänge, in Leinwand gebunden. 96 Seiten 10½×8. Preis: 30 Cts. und 8 Cts. Porto.

Diese „Kirchenglocke“ erscheint schon seit 1897 als ein Tractatblatt, das die Leute zum fleißigen und andächtigen Kirchenbesuch ermuntern will. Dazu ist der mannigfache Inhalt der einzelnen Nummern auch wohl geeignet. Aber auch die vorliegenden gebundenen sechs Jahrgänge lassen sich noch verwerthen, z.\,B. im Confirmandenunterricht. Der Herausgeber schreibt: „Wenn ich meinen Confirmanden sage: Wenn ihr heute gut aufmerkt und gut antwortet, dann will ich euch noch etwas vorlesen, habe ich eine ruhige, aufmerksame Klasse.“ Und es sind in der That gar manche Sachen in den 24 Nummern, für die sich Kinder besonders interessiren werden. L. F.

\section[Experimental Religion]{Experimental Religion}
The Experiences of Christopher from his Awaking to his Falling Asleep in Jesus. By F. K.: VIRGINIUS. 198 Seiten 7½×5, in Leinwand gebunden. Augusta Publishing Company, Crimora, Va. Auch zu beziehen vom Concordia Publishing House, St. Louis, Mo. Preis: 75 Cts.

Dieses schon vor einigen Monaten erschienene Buch eines unter uns wohlbekannten Verfassers können wir nur empfehlen. Es schildert das Leben eines jungen, weltlichgesinnten Mannes, der in einem Gottesdienste zur Erkenntniß seiner Sünde kommt und den Weg zum Leben sucht. Erst wird er Irrwege geführt, auf denen er nicht zur Ruhe gelangt, hört dann aber in einer Osterpredigt die Lehre von der Rechtfertigung aus Gnaden durch den Glauben und kommt dadurch zur Gewißheit der Vergebung seiner Sünden und zum Frieden mit Gott. Anfechtungen bleiben freilich nicht aus, die ihm aber zum Besten dienen und unter denen er immer tiefer in die Schrift und in die Erkenntniß der heilsamen Lehre von seinem Seelsorger eingeführt wird. Es kommt auch eine Zeit, da er geistlich sicher wird und in Sünde fällt, durch Gottes Gnade jedoch wieder aufsteht und dann trotz mancher Heimsuchungen ein glückliches und segensreiches Leben bis zu seinem seligen Tod führt. Dies alles ist so schlicht und doch so interessant beschrieben und zeigt in der Form einer Erzählung so schön und klar den Weg zum Heile, daß wir dieses Buch namentlich auch in englischen Kreisen recht verbreitet sehen möchten. Es würde da gewiß von manchem gelesen, der einen Tractat oder eine Lehrabhandlung nicht beachten würde. Man hat das Büchlein nicht unpassend das lutherische „Pilgrim’s Progress“ genannt. Während aber das bekannte Werk John Bunyans nicht frei ist vom Sauerteig der reformirten Irrlehre, ist hier alles dem Vorbilde der heilsamen Worte gemäß. L. F.

\section[Food on the Way]{Food on the Way}
Pastoral Letters, Prayers, and Hymns for our Youth newly Confirmed, submitted by C. M. Zorn. Rendered from the German by A. W. M. St. Louis, Mo. Published by F. Dette. 1904. 127 Seiten 5½×4. Auch zu beziehen vom Concordia Publishing House, St. Louis, Mo. Preis gebunden: 20 Cts., das Dutzend \$1.80; mit Goldschnitt 30 Cts., das Dutzend \$3.00.

Dies ist eine gute von einem bekannten Pastor aus unserer englischen Schwestersynode besorgte Uebersetzung des von uns schon zweimal empfohlenen passenden Confirmandenbüchleins P. Zorns. (Vergl. „Lutheraner“, Jahrg. 58, S. 58. Jahrg. 59, S. 74.) Möge es auch in der englischen Sprache von vielen Confirmanden fleißig gebraucht werden und ihnen zum Segen gereichen für Zeit und Ewigkeit. L. F.

\section{Christus lebt!}
Osterlied für Männerchor von H. B. Bröhl, 3606 S. Hermitage Ave., Chicago, Ill. 3 Seiten 11×7½. Preis: 15 Cts., das Dutzend \$1.50 portofrei.
